% !PRINTED 204
\weberpage{204}
%<<CONT
\weq{4}{\begin{gathered}
H + 4\,a_2'\,f = -3\,\psi_1^2\\
H + 4\,a_2''\,f = -3\,\psi_2^2\\
H + 4\,a_2'''\,f = -3\,\psi_3^2.
\end{gathered}}

Of the three functions $\psi_1, \psi_2, \psi_3$ \gesperrt{no two can have
a common divisor}, provided the discriminant
$D$ is assumed different from zero. For
then the three quantities $a_2', a_2'', a_2'''$ are also different from one another.
If then $\psi_1$ and $\psi_2$ vanish, $H$ and
$f$ must vanish, i.\,e.\ a common divisor of $\psi_1$ and $\psi_2$
would have to be a common divisor of $H$ and $f$. But now $\xi$
was an arbitrary linear divisor of $f$. By (2) this can be
a divisor of $H$ only if $a_3' = 0$, that is, again, only if
the discriminant $D$ vanishes [§.\,65, (11)].

Now the functional determinant $T$ of $f$ and $H$ is identical
with the functional determinant of $f$ and $H + \lambda f$, whatever
$\lambda$ may be, and from this it follows, when $\lambda = 4\,a_2'$ is
set,
\[
T = -\tfrac{1}{2}\,\psi_1\,[F'(\xi)\,\psi_1'(\eta) - F'(\eta)\,\psi_1'(\xi)].
\]

Thus $T$ is divisible by $\psi_1$, and for the same reasons
by $\psi_2$ and $\psi_3$, and consequently $T$ is, up to a factor independent
of $\xi, \eta$, identical with the product $\psi_1\,\psi_2\,\psi_3$.
If we accordingly form the product of the three equations (4),
it follows, having regard to the cubic equation §.\,65, (12),
which the quantities $a_2', a_2'', a_2'''$ satisfy:
\[
H^3 - 48\,A\,H\,f^2 - 64\,B\,f^3 = c\,T^2,
\]
where $c$ is a constant still to be determined. By §.\,64, (12)
$c$ remains unchanged under a linear transformation, and one
therefore finds its value by setting the terms with the highest
power of $\xi$ on the two sides of the normal form (1), (2)
equal to each other, $c = -27$.

We thus have between the invariants and covariants
the following identical relation
\weq{5}{H^3 - 48\,A\,H\,f^2 - 64\,B\,f^3 = -27\,T^2.}

The functions $\psi_1, \psi_2, \psi_3$ are likewise covariants, though
with coefficients which are not rational in the coefficients
of $f$. We can easily express them in terms of the roots $\alpha, \beta, \gamma, \delta$
of $f = 0$. If for simplification one sets $y = 1$,
then by §.\,65, (3) and (5) it follows that
%CONT>>
